\section{Problem Drift} \label{sec:definition}

%To study the effect of MAD, especially concerning the convergence of task performance,
We introduce the concept of \textbf{problem drift} in MAD, which we define as the systematic decay in turn-based task performance as agents gradually diverge from the initial task. 
To quantify this effect, we propose \metric\ to measure the strength of the decay over a number of consecutive turns.

\paragraph{\textbf{\metric.}}
We describe how the quality of a solution differs in one turn from the previous turn.
Let $(x, y)$ be a pair of input and gold label for a task example in a dataset to be solved. 
A multi-agent system discusses for $N$ rounds.
At the end of each round $r \in (1, ..., N)$, the agents try to solve $x$ by providing an answer $\hat{y}^{(r)}$. 
%Let $\hat{y}_r \in A_r$ be the ground-truth solution at round $r$ and  
Let $M(\hat{y}^{(r)}, y) \in [0,1]$ be a downstream performance metric that compares the ground-truth solution to the solution produced by the MAD, where $1$ corresponds to the solution and $0$ corresponds to a wrong solution.
Values between $0$ and $1$ can indicate partial solutions for non-binary solutions.
We define \textit{problem focus}, further called \metric, as:
\vspace*{2mm}
\begin{equation} \label{def:focus}
\small
FOCUS_r = M(\hat{y}^{(r)}, y) - M(\hat{y}^{(r-1)}, y) \in [-1,1]
\end{equation}

\noindent
If $FOCUS_r < 0$, the solution's quality degrades during round $r$. If $FOCUS_r > 0$, the solution improves.
\metric\ purposely depends on task-performance as an outcome-based score.
We also test a semantic-based alternative in \Cref{app:alternative_metric} but find two main limitations that make it unattractive for this work.

\paragraph{\textbf{Problem Drift.}}
We define that an ongoing discussion has \textit{problem drift} if the current solution obtains lower performance than the solution of a previous debate round over multiple rounds.
Let $M$ be the number of consecutive discussion rounds for one task example.
A debate has problem drift from round $1$ to $M$ of strength $\theta \in [-1,0]$ if:
\begin{equation} \label{def:drift}
\small
FOCUS_{1,M} = \sum_{r=1}^{M} FOCUS_r = \theta
\end{equation}
We define problem drift by the simple sum over \metric\ values to model the shift in task performance across debate rounds.
% An alternative approach could also favor more recent rounds. 
% However, this would penalize temporary explorations that can be later recovered. 
% In contrast, a momentum term would not accurately capture the severity of problem drift at specific rounds.

\paragraph{\textbf{Recovery.}}
We say an example recovered from problem drift when the discussion gets back to or improves upon the performance before the problem drift occurred.
We define a discussion as recovered from problem drift if $\exists N,M \text{ where } N > M \text{ s.t.}$:

% \begin{equation} \label{def:recovery}
% R_n = \frac{\sum_{r=i}^n \mathbf{1}(\sum_{j=1}^{M} FOCUS_j >= 0)}{\sum_{r=i}^n \mathbf{1}(\sum_{j=1}^{N} FOCUS_j< 0)} \times 100, \quad \forall r \in [i, n]
% \end{equation}

% \begin{equation} \label{def:recovery}
%     \small
%     R_P = \frac{\sum_{d=1}^P \mathbb{1} (\sum_{r=1}^{N} FOCUS_r >= 0)}{\sum_{d=1}^P \mathbb{1} (\exists M<N \text{ s.t. } \sum_{r=1}^{M} FOCUS_r < 0)}
% \end{equation}

\vspace*{-4mm}
\begin{equation} \label{def:recovery}
    \small
    \sum_{r=1}^{N} FOCUS_r >= 0 \quad \wedge \quad \sum_{r=1}^{M} FOCUS_r < 0
\end{equation}